%% !TEX program = xelatex
\documentclass[12pt,final]{novelette} % font size, twoside, twocolumn, draft-final, use extarticle for 8-20pt

\title{Diff Eq Midterm 2544}
\author{Kaka Cring}
\date{\today}

\begin{document}

\makeheading{ตัวอย่างข้อสอบ 2301312 สมการเชิงอนุพันธ์ ปรับปรุงจากข้อสอบกลางภาคปลาย ปีการศึกษา 2544}

\subsubsection*{PART A}

\begin{question}
    \item {}
    จงหารอนสเกียนของ $e^x, \cos x$ และ $\sin x$ \\
    และจากรอนสเกียนที่หาได้จงบอกว่าฟังก์ชันทั้งสามเป็นอิสระเชิงเส้นต่อกันหรือไม่บนช่วง $(-\infty, \infty)$ เพราะเหตุใด \qmark{3}

    % ----------------------------------------
    \begin{align*}
        W(e^x, \cos x, \sin x) & =
        \begin{vmatrix}
            e^x & \cos x  & \sin x  \\
            e^x & -\sin x & \cos x  \\
            e^x & -\cos x & -\sin x
        \end{vmatrix}                           \\[1ex]
                               & =
        e^x\Bigg(
        \begin{vmatrix}
                -\sin x & \cos x  \\
                -\cos x & -\sin x
            \end{vmatrix}                                 \\
                               & \qquad -
        \begin{vmatrix}
                \cos x  & \sin x  \\
                -\cos x & -\sin x
            \end{vmatrix}                                 \\
                               & \qquad +
        \begin{vmatrix}
                \cos x  & \sin x \\
                -\sin x & \cos x
            \end{vmatrix}
        \Bigg)                                            \\
                               & =        e^x (1 - 0 + 1) \\
                               & = \boxed{2e^x}
    \end{align*}

    Since Wronskian is not equal to $0$ anywhere, then all three functions are linearly independent on $(-\infty, \infty)$.
    % ----------------------------------------

    \item {}
    จงให้เหตุผลว่าทําไมสมการ $(1-x)y'' + xy' - y = 2x^3 - 6x^2 + 6x$ จึงเป็นปรกติบนช่วง $(1, \infty)$
    และจงหาผลเฉลยบริบูรณ์ของสมการนี้บนช่วง $(1, \infty)$ โดยให้พิจารณาจากฟังก์ชัน $e^x, x$ และ $x^3$ \qmark{8}

    % ----------------------------------------
    \bigskip
    Since every coefficients are polynomials of $x$, they are differentiable everywhere.
    The condition that $1-x \neq 0$ is satisfied for any $x \in (1, \infty)$.
    Therefore this equation is normal on $(1, \infty)$.

    \medskip
    Solve the relative homogenous equation. \\
    Try $y = c_1 e^x$;
    $$
        (1-x)c_1 e^x + x c_1 e^x - c_1 e^x = 0
    $$
    is true for all $c_1$, so $c_1 e^x$ is a solution.

    Try $y = c_2 x$;
    $$
        x c_2 - c_2 x = 0
    $$
    is true for all $c_2$, so $c_2 x$ is a solution.

    \medskip
    Since these 2 solutions are linearly independent, the complementary function is
    $$
        y_c = c_1 e^x + c_2 x
    $$

    \medskip
    Solve the particular integral. \\
    Try $y = x^3$
    \begin{align*}
        6 (1-x) x + 3 x x^2 - x^3 & = 2x^3 - 6x^2 + 6x \\
        6x - 6x^2 + 2x^3          & = 2x^3 - 6x^2 + 6x \\
    \end{align*}
    Therefore $x^3$ is the particular integral.

    \medskip
    The complete solution is
    $$
        \boxed{y = c_1 e^x + c_2 x + x^3}
    $$
    % ----------------------------------------

    \item {}
    ลวดสปริงเส้นหนึ่งถูกยึดปลายบนติดกับเพดาน
    เมื่อนําวัตถุหนัก $6$ ปอนด์ มาผูกติดไว้ที่ปลายล่างของลวดสปริงจะทําให้ขดลวดสปริงยืดออก $6$ นิ้ว
    ถ้าดึงวัตถุนี้ลงมาต่ำกว่าตำแหน่งสมดุล $4$ นิ้ว แล้วปล่อยด้วยการผลักวัตถุขึ้นด้วยความเร็ว $\frac83$ ฟุต/วินาที \\
    จงหาสมการของการเคลื่อนที่ของวัตถุพร้อมทั้งบอกค่าแอมพลิจูด และมุมทิ้งช่วง \qmark{10}

    % ----------------------------------------
    \bigskip
    The governing equation is (positive $x$ direction is downward)
    $$
        mg - kx = m\ddot{x}
    $$
    Using the equilibrium case
    \begin{align*}
        mg & = kx_0                                                \\
        k  & = \frac{mg}{x_0}                                      \\
        k  & = \frac{\SI{6}{lb}\cdot\SI{32}{ft/s^2}}{\SI{0.5}{ft}} \\
        k  & = \SI{384}{lb/s^2}
    \end{align*}
    Rewrite the equation as
    \begin{align*}
        k(x_0 - x)                           & = m\ddot{x}                    \\
        (\SI{384}{lb/s^2})(\SI{0.5}{ft} - x) & = (\SI{6}{lb})\ddot{x}         \\
        32 - 64 x                            & = \ddot{x}                     \\
        (D^2+8^2)x                           & = 32                           \\
        x(t)                                 & = A \sin(8t + \phi_0)\,\si{ft}
    \end{align*}
    Use initial conditions
    \begin{align*}
        x_0               & = A \sin(\phi_0)\,\si{ft}    \\
        \frac13           & = A \sin(\phi_0) \tag{i}     \\
        x'(0)             & = 8A \cos(\phi_0)\,\si{ft/s} \\
        -\frac83          & = 8A \cos(\phi_0)            \\
        -\frac13          & = A \cos(\phi_0) \tag{ii}    \\
        \intertext{$(\mathrm i)^2 + (\mathrm{ii})^2$;}
        \frac19 + \frac19 & = A^2                        \\
        A                 & = \frac{\sqrt2}{3} \tag{iii} \\
        \intertext{Substitute (iii) into (i);}
        \sin(\phi_0)      & = \frac1{\sqrt2}             \\
        \phi_0            & = \frac\pi4, \pi - \frac\pi4 \\
        \intertext{Since $\frac\pi2 < \phi_0 < \pi$;}
        \phi_0            & = \frac{3\pi}{4}
    \end{align*}
    The motion equation is
    $$
        \boxed{x(t) = \frac{\sqrt2}{3}\sin\left(8t+\frac{3\pi}{4}\right)\,\si{ft}}
    $$
    Amplitude is $\boxed{\frac{\sqrt2}3\,\si{ft}}$ and angle of departure is $\boxed{\frac{3\pi}{4}\,\si{rad}}$.
    % ----------------------------------------

    \item {}
    สมการของการเคลื่อนที่ของวัตถุพร้อมทั้งบอกค่าแอมพลิจูด และมุมทิ้งช่วง \\
    ถ้าเราทราบว่าสมการของการเคลื่อนที่ของวัตถุคือ $x(t) = \frac13 \sin(\pi t - \frac\pi4)$
    \begin{question}
        \item
        จงหา แอมพลิจูด คาบ และความถี่ของการเคลื่อนที่ \qmark{1.5}

        % ----------------------------------------
        \bigskip
        Amplitude is $\boxed{\frac13}$. Period is $\boxed{2}$. Frequency is $\boxed{\frac12}$.
        % ----------------------------------------

        \item
        จงหาตําแหน่ง ความเร็ว ความเร่งของการเคลื่อนที่ เมื่อเวลา $1$ วินาที หลังจากปล่อยให้วัตถุเกิดการเคลื่อนที่
        พร้อมทั้งอธิบายลักษณะของการเคลื่อนที่ของวัตถุ ณ ขณะนั้น \qmark{3.5}

        % ----------------------------------------
        \phantom{}
        \vspace*{-1em}
        \begin{align*}
            x(1)   & = \frac13 \sin\left(\pi - \frac\pi4\right)   \\
                   & = \boxed{\frac1{3\sqrt2}}                    \\
            x'(1)  & = \frac\pi3 \cos\left(\pi - \frac\pi4\right) \\
                   & = \boxed{\frac{\pi}{3\sqrt2}}                \\
            x''(1) & = \boxed{-\frac{\pi^2}{3\sqrt2}}
        \end{align*}

        The object is travelling in $+x$ direction and is getting accelerated into $-x$ direction.
        % ----------------------------------------
    \end{question}
\end{question}

\subsubsection*{PART B}

\begin{question}[resume]
    \item {}
    จงเติมเฉพาะคำตอบลงในช่องว่างที่กำหนดให้โดยไม่ต้องแสดงวิธีทำ \qmark{20}
    \begin{question}
        \item
        ผลเฉลยบริบูรณ์ของสมการ $D(D-3)y = 0$ \qmark{2}

        % ----------------------------------------
        $$ \boxed{y = c_1 + c_2 e^{3x}} $$
        % ----------------------------------------

        \item
        ผลเฉลยบริบูรณ์ของสมการ $(D^2 - 2D + 2)^2 y = 0$ \qmark{2}

        % ----------------------------------------
        $$ \boxed{e^x (c_1 + c_2 x + c_3 x^2 + c_4 x^3)} $$
        % ----------------------------------------

        \pagebreak

        \item
        ผลเฉลยบริบูรณ์ของสมการสมการเอกพันธ์ุสัมพัทธ์ของสมการ $(D-2)^2(D^2+2D+5)y = e^x$ \qmark{2}

        % ----------------------------------------
        $$ \boxed{e^{2x} (c_1 + c_2 x) + e^{-x}(c_3 \cos(2x) + c_4 \sin(2x))} $$
        % ----------------------------------------

        \item
        ผลเฉลยเฉพาะหรือปริพันธ์เฉพาะของสมการ $(D-2)^2(D^2+2D+5)y = e^x$ \qmark{2}

        % ----------------------------------------
        $$ \boxed{\frac18 e^x} $$
        % ----------------------------------------

        \item
        ผลเฉลยบริบูรณ์ของสมการสมการเอกพันธ์ุสัมพัทธ์ของสมการ $y^{(4)} + 4y'' = \sin x$ \qmark{2}

        % ----------------------------------------
        $$ \boxed{c_1 + c_2 x + c_3 \cos(2x) + c_4 \sin(2x)} $$
        % ----------------------------------------

        \item
        ผลเฉลยเฉพาะหรือปริพันธ์เฉพาะของสมการ $y^{(4)} + 4y'' = \sin x$ \qmark{2}

        % ----------------------------------------
        $$ \boxed{-\frac{1}{3} \sin x} $$
        % ----------------------------------------  

        \item
        ตัวดําเนินการเชิงอนุพันธ์อันดับต่ำที่สุดที่ลบล้าง $xe^{2x}$ \qmark{2}

        % ----------------------------------------
        $$ \boxed{(D-2)^2} $$
        % ----------------------------------------

        \item
        ตัวดําเนินการเชิงอนุพันธ์อันดับต่ำที่สุดที่ลบล้าง $e^{-x}\sin 3x$ \qmark{2}

        % ----------------------------------------
        $$ \boxed{(D-1)^2 + 9} $$
        % ----------------------------------------

        \item
        สมการเชิงอนุพันธ์สามัญเชิงเส้นเอกพันธุ์ที่มีสัมประสิทธิ์เป็นค่าคงตัวอันดับต่ำที่สุดที่มี $e^x$ และ $x^2$ เป็นผลเฉลย \qmark{2}

        % ----------------------------------------
        $$ \boxed{(D-1)D^3} $$
        % ----------------------------------------

        \pagebreak

        \item
        สมการเชิงอนุพันธ์สามัญเชิงเส้นเอกพันธุ์ที่มีสัมประสิทธิ์เป็นค่าคงตัวอันดับต่ำที่สุดที่มี $x \cos 2x$ และ $x^2 e^x$ เป็นผลเฉลย \qmark{2}

        % ----------------------------------------
        $$ \boxed{(D^2 + 4)^2 (D-1)^3} $$
        % ----------------------------------------
    \end{question}

    \item {}
    จงหาผลเฉลยเฉพาะของสมการ
    $y''' + y'' = 0; \quad y(0) = 2, \quad y'(0) = 0, \quad y''(0) = 1$ \qmark{5}

    % ----------------------------------------
    \bigskip
    Rewrite the equation as
    $$ D^2(D - 1) y = 0 $$
    The general solution is
    $$ y = c_1 + c_2 x + c_3 e^x $$
    Solve for the coefficients
    \begin{align*}
        y''(0) & = c_3 e^0                 \\
        1      & = c_3                     \\
        y'(0)  & = c_2 + c_3 e^0           \\
        0      & = c_2 + 1                 \\
        -1     & = c_2                     \\
        y(0)   & = c_1 + c_2 (0) + c_3 e^0 \\
        2      & = c_1 + 1                 \\
        1      & = c_1
    \end{align*}
    Therefore the particular solution is
    $$ \boxed{y = 1 - x + e^x} $$
    % ----------------------------------------
\end{question}

\subsubsection*{PART C}

\begin{question}[resume]
    \item {}
    จงหาปริพันธ์เฉพาะของสมการ $y'' - y = \frac32(1-\cos2x)$ โดยวิธีเทียบสัมประสิทธิ์ \qmark{7}

    % ----------------------------------------
    \bigskip
    Find complementary function
    \begin{align*}
        (D^2-1) y & = 0                      \\
        y_c       & = c_1 e^{x} + c_2 e^{-x}
    \end{align*}

    Find annihilator differential operator of $\frac32 (1 - \cos 2x)$
    $$
        Q(D) = D(D^2 + 4)
    $$

    Solve general solution of this homogenous equation
    \begin{align*}
        Q(D)P(D) y        & = 0                                                      \\
        D(D^2+4)(D^2-1) y & = 0                                                      \\
        y                 & = c_1 + c_2 \sin 2x + c_3 \cos 2x + c_4 e^x + c_5 e^{-x} \\
        y_q               & = c_1 + c_2 \sin 2x + c_3 \cos 2x
    \end{align*}

    Suppose the particular integral is in the same form as $y_q$
    \begin{align*}
        y_p                                 & = k_1 + k_2 \sin 2x + k_3 \cos 2x \\
        (D^2-1)y_p                          & = \frac32 (1 - \cos 2x)           \\
        -k_1 - 5 k_2\sin 2x - 5 k_3 \cos 2x & = \frac32 (1 - \cos 2x)           \\
        k_1                                 & = -\frac{3}{2}                    \\
        k_2                                 & = 0                               \\
        k_3                                 & = \frac{3}{10}
    \end{align*}

    The particular integral is
    $$
        \boxed{y_p = - \frac{3}{2} + \frac{3}{10} \cos 2x}
    $$
    % ----------------------------------------

    \pagebreak

    \item {}
    จงใช้ตัวดำเนินการผกผันหาปริพันธ์เฉพาะของสมการ $(D^2 - 9D + 18)y = e^{e^{-3x}}$ \qmark{5}

    % ----------------------------------------
    \vspace*{-1em}
    \begin{align*}
        (D-6)(D-3)y & = e^{e^{-3x}}                                                     \\
        y           & = \frac{1}{(D-6)(D-3)} e^{e^{-3x}}                                \\
                    & = \frac{1}{D-6} e^{3x} \int e^{-3x} e^{e^{-3x}} \dd{x}            \\
                    & = \frac{1}{D-6} e^{3x} \int -\frac{1}{3} e^{e^{-3x}} \dd(e^{-3x}) \\
                    & = \frac{1}{D-6} e^{3x} -\frac{1}{3} e^{e^{-3x}}                   \\
                    & = -\frac13 e^{6x} \int e^{-6x} e^{3x} e^{e^{-3x}} \dd{x}          \\
                    & = -\frac13 e^{6x} \int -\frac{1}{3} e^{e^{-3x}} \dd(e^{-3x})      \\
                    & = \boxed{\frac19 e^{6x+e^{-3x}}}
    \end{align*}
    % ----------------------------------------

    \item {}
    จงใช้ตัวดำเนินการผกผันหาปริพันธ์เฉพาะของสมการ $(D^2+9)y = x \cos x$ \qmark{6}

    % ----------------------------------------
    \vspace*{-1em}
    \begin{align*}
        \frac{1}{D^2+9} x \cos x & = \frac{1}{D^2+9} x \Re\left(e^{ix}\right)                                \\
                                 & = \Re\left(\frac{1}{D^2+9} x e^{ix}\right)                                \\
                                 & = \Re\left(e^{ix} \frac{1}{(D+i)^2+9} x\right)                            \\
                                 & = \Re\left(e^{ix} \frac{1}{D^2+2iD-1+9} x\right)                          \\
                                 & = \Re\left(\frac{e^{ix}}8 \frac{1}{1-(-D^2-2iD)/8} x\right)               \\
                                 & = \Re\left(\frac{e^{ix}}8 \left(1-\frac{2iD}{8}\right) x\right)           \\
                                 & = \Re\left(\frac{e^{ix}}8 \left(x - \frac{i}{4}\right)\right)             \\
                                 & = \frac18 \Re\left((\cos x + i \sin x)\left(x - \frac{i}{4}\right)\right) \\
        \Aboxed{y_p              & = \frac18 x \cos x + \frac1{32} \sin x}
    \end{align*}
    % ----------------------------------------

    \pagebreak

    \item {}
    จงหาปริพันธ์เฉพาะของสมการ $(D^3 + D^2 - 2D)y = e^{2x}$ โดยวิธีแปรพารามิเตอร์ \qmark{7}

    % ----------------------------------------
    \bigskip
    Find complementary function
    \begin{align*}
        m^3 + m^2 - 2m & = 0                             \\
        m^2 + m - 2    & = 0 \quad \text{or} \quad m = 0 \\
        (m-1)(m+2)     & = 0                             \\
        m = -2, 1
    \end{align*}
    Therefore
    $$
        y_c = c_1 + c_2 e^{-2x} + c_3 e^x
    $$

    Solve for Wronskian
    \begin{align*}
        W(1,e^{-2x},e^x) & =
        \begin{vmatrix}
            -2e^{-2x} & e^x \\
            4e^{-2x}  & e^x
        \end{vmatrix}                \\
                         & = -6 e^{-x}
    \end{align*}
    Let
    $$
        y_p = u_1(x) + u_2(x) e^{-2x} + u_3(x) e^x
    $$

    Find each component
    \begin{align*}
        \intertext{Find $u_1$;}
        u'_1 & = e^{2x}\frac{
            \begin{vmatrix}
                e^{-2x}   & e^x \\
                -2e^{-2x} & e^x
            \end{vmatrix}
        }{W}                                      \\
             & = e^{2x}\frac{2e^{-x}}{-6e^{-x}}   \\
             & = -\frac13e^{2x}                   \\
        u_1  & = - \frac16 e^{2x}                 \\
        \intertext{Find $u_2$;}
        u'_2 & = - e^{2x}\frac{
            \begin{vmatrix}
                1 & e^x \\
                0 & e^x
            \end{vmatrix}
        }{W}                                      \\
             & = e^{2x}\frac{e^x}{-6e^{-x}}       \\
             & = -\frac16 e^{4x}                  \\
        u_2  & = - \frac{1}{24} e^{2x}            \\
        \intertext{Find $u_3$;}
        u'_3 & = e^{2x}\frac{
            \begin{vmatrix}
                1 & e^{-2x}   \\
                0 & -2e^{-2x}
            \end{vmatrix}
        }{W}                                      \\
             & = e^{2x}\frac{-2e^{-2x}}{-6e^{-x}} \\
             & = \frac13 e^x                      \\
        u_3  & = \frac13 e^x
    \end{align*}
    The particular integral is
    \begin{align*}
        y_p         & = -\frac16e^{2x} - \frac1{24} e^{2x} + \frac13 e^{2x} \\
        \Aboxed{y_p & = \frac18 e^{2x}}
    \end{align*}
    % ----------------------------------------
\end{question}

\subsubsection*{PART D}

\begin{question}[resume]
    \item {}
    จงหาปริพันธ์เฉพาะของสมการ $(D^2 - 3D + 1)y = x^2 + 15 \sin 2x$ \qmark{8}

    % ----------------------------------------
    \bigskip
    First part;
    \begin{align*}
        \frac{1}{D^2-3D+1}x^2 & = (1 - (D^2 - 3D) + (D^2 - 3D)^2) x^2 \\
                              & = x^2 - 2 + 6x + 18                   \\
                              & = x^2 + 6x + 16
    \end{align*}
    Second part;
    \begin{align*}
        \frac{1}{D^2-3D+1}\sin2x & = \frac{1}{-2^2-3D+1}\sin 2x                                                          \\
                                 & = -\frac13 \frac{1}{D+1}\sin 2x                                                       \\
                                 & = -\frac13 e^{-x} \int e^x \sin 2x \dd{x}                                             \\
                                 & = -\frac13 e^{-x} \frac45 \left(\frac12 e^x (-\cos 2x) + \frac14 e^x (\sin 2x)\right) \\
                                 & = \frac{1}{15} (2 \cos 2x - \sin 2x)
    \end{align*}
    The particular integral is
    $$
        \boxed{y_p = x^2 + 6x + 16 + 2 \cos 2x - \sin 2x}
    $$
    % ----------------------------------------

    \item {}
    จงหาผลเฉลยบริบูรณ์ของสมการ $4x^2y'' + 8xy' + y = x^3$ \qmark{6}

    % ----------------------------------------
    \bigskip
    This is a Cauchy-Euler equation.
    Substitute $x = e^z$, then
    $$ x^n \dv[n]{y}{x} = D(D-1)\cdots(D-n+1)y \ \text{if}\  D = \dv{z} $$

    The equation becomes
    \begin{align*}
        4 D(D-1)y + 8Dy + y          & = 0                          \\
        (4D^2 + 4D + 1) y            & = 0                          \\
        \left(D + \frac12\right)^2 y & = 0                          \\
        y_c                          & = e^{-z/2} (c_1 + c_2 z)     \\
        y_c                          & = x^{-1/2} (c_1 + c_2 \ln x)
    \end{align*}

    The particular integral is
    \begin{align*}
        y_p & = \frac{1}{4(D + 1/2)^2} e^{3z}                    \\
            & = \frac14 e^{-z/2} \iint e^{z/2} e^{3z} (\dd{z})^2 \\
            & = \frac1{14} e^{-z/2} \int e^{7z/2} \dd{z}         \\
            & = \frac1{49} e^{3z}                                \\
            & = \frac1{49} x^3
    \end{align*}

    The complete solution is
    $$
        \boxed{y = x^{-1/2} (c_1 + c_2 \ln x) + \frac1{49} x^3}
    $$

    % ----------------------------------------

    \item {}
    จงหาผลเฉลยในรูปอนุกรมกำลังรอบจุดสามัญ $x=0$ ของสมการ $y'' - xy' - y = 0$ \\
    (เขียนผลเฉลยให้เห็นอย่างน้อยถึงพจน์ $x^5$) \qmark{10}

    % ----------------------------------------
    \bigskip
    Let $ \displaystyle y = \sum_{n=0}^\infty a_n x^n $.
    \begin{align*}
        \sum_{n=2}^\infty n(n-1) a_n x^{n-2} - x \sum_{n=1}^\infty n a_n x^{n-1} - \sum_{n=0}^\infty a_n x^n & = 0 \\
        \sum_{n=0}^\infty (n+2)(n+1) a_{n+2} x^n - \sum_{n=1}^\infty n a_n x^n - \sum_{n=0}^\infty a_n x^n   & = 0 \\
        2 a_2 - a_0 + \sum_{n=1}^\infty (n+2)(n+1) a_{n+2} x^n - n a_n x^n - a_n x^n                         & = 0
    \end{align*}

    Solve for coefficients
    \begin{align*}
        a_2     & = \frac{a_0}{2}     \\
        a_{n+2} & = \frac{1}{n+2} a_n \\
    \end{align*}

    Let $a_0 = c_1$ and $a_1 = c_2$, then
    \begin{align*}
        a_2 & = \frac{c_1}{2}  \\
        a_3 & = \frac{c_2}{3}  \\
        a_4 & = \frac{c_1}{8}  \\
        a_5 & = \frac{c_2}{15}
    \end{align*}

    The solution is
    $$
        \boxed{y = c_1\left(1 + \frac{x^2}{2} + \frac{x^4}{8} + \cdots\right) + c_2\left(x + \frac{x^3}{3} + \frac{x^5}{15} + \cdots\right)}
    $$
    % ----------------------------------------
\end{question}

\end{document}